

\documentclass{article}
\usepackage[utf8]{inputenc}
\usepackage{authblk}
\usepackage{setspace}
\usepackage[margin=1.25in]{geometry}
\usepackage{graphicx}
\graphicspath{ {./figures/} }
\usepackage{subcaption}
\usepackage{amsmath}
\usepackage{lineno}
\linenumbers

%%%%%% Bibliography %%%%%%
% Replace "sample" in the \addbibresource line below with the name of your .bib file.
\usepackage[style=nejm, 
citestyle=numeric-comp,
sorting=none]{biblatex}
\addbibresource{sample.bib}

%%%%%% Title %%%%%%
% Full titles can be a maximum of 100 characters, including spaces. 
% Title Format: Use title case, capitalizing the first letter of each word, except for certain small words, such as articles and short prepositions
\title{Health Data Science \LaTeX\ Template}

%%%%%% Authors %%%%%%
% Authors should be listed in order of contribution to the paper, by first name, then middle initial (if any), followed by last name.
% Authors should be listed in the order in which they will appear in the published version if the manuscript is accepted. 
% Use an asterisk (*) to identify the corresponding author, and be sure to include that person’s e-mail address. Use symbols (in this order: †, ‡, §, ||, ¶, #, ††, ‡‡, etc.) for author notes, such as present addresses, “These authors contributed equally to this work” notations, and similar information.
% You can include group authors, but please include a list of the actual authors (the group members) in the Supplementary Materials.
\author[1*$\dag$]{Author One}
\author[2$\dag$]{Author Two}
\author[2]{Author Three}
\author[1,2]{Author Four}

%%%%%% Affiliations %%%%%%
\affil[1]{Department of Physics, A University, City, Country.}
\affil[2]{Department of Astronomy, B University, City, Country.}
\affil[*]{Address correspondence to: email@email.com}
\affil[$\dag$]{These authors contributed equally to this work.}

%%%%%% Date %%%%%%
% Date is optional
\date{}

%%%%%% Spacing %%%%%%
% Use paragraph spacing of 1.5 or 2 (for double spacing, use command \doublespacing)
\onehalfspacing

\begin{document}

\maketitle

%%%%%% Abstract %%%%%%
\begin{abstract}
The abstract should be a single paragraph written in plain language that a general reader can understand. Do not include citations, figures, tables, or undefined abbreviations in the abstract. Any abbreviations that appear in the title should be defined in the abstract. The length should be 200 words and not exceed 250 words, to include: 
\begin{itemize}
    \item An opening sentence that states the question/problem addressed by the research AND
    \item Enough background content to give context to the study AND
    \item A brief statement of primary results AND
    \item A short concluding sentence.
\end{itemize} 
\end{abstract}

%%%%% Citations in the text %%%%%%
\subsection*{Citations}
Citations of references in the text should be identified using numbers in square brackets e.g., ``as discussed by Cui \cite{Cui1}'' or ``as discussed elsewhere \cite{Cui1,Ninomiya1,Li1,Wang1,Yang1}.'' All references should be cited within the text and uncited references will be removed. 

As an example, this template includes a ``sample.bib'' file containing the references in BibTeX.

%%%%%% Equations %%%%%%
\subsection*{Equations}
Equations should be provided in a text format, rather than as an image. Equations should be numbered consecutively, in round brackets, on the right-hand side of the page by using the ``\textbackslash begin\{equation\}'' command. They should be referred to as Equation 1, etc. in the main text.

\medskip For example, see Equation \ref{eq:1} and Equation \ref{eq:2} below.
\begin{equation} \label{eq:1}
    a^2 + b^2 = c^2
\end{equation}
\begin{equation} \label{eq:2}
\begin{split}
A & = \frac{\pi r^2}{2} \\
 & = \frac{1}{2} \pi r^2
\end{split}
\end{equation}

%%%%%% Figures and Tables %%%%%%
\subsection*{Figures and Tables}
For initial submissions, Figures and Tables should be embedded within the main submission file at their first mention in the text. This will facilitate evaluation of the paper. Figures should be called out within the text and numbered in the order of their citation (with title and caption formatted as below). See below for detailed instructions on preparation of and preferred formats for your figures.

Tables should supplement, not duplicate, the text. They should be called out within the text and numbered in the order of their citation in the text, as noted in Table \ref{tab:1}. As with Figures, please embed Tables within the text of the initial submission as near to their first reference as possible. Each table should be preceded by a legend that starts with a short boldface title beginning with (for example) \textbf{Table 1}. Every vertical column should have a heading, followed by a unit of measure (if any) in parentheses. Units should not change within a column. Centered headings of the body of the table can be used to break the entries into groups. Footnotes should contain information relevant to specific cells of the table; use lowercase letters in alphabetical order, as needed: a, b, c, etc. (Don’t use footnotes in column heads; include any such details in sentence form in the table legend.)

\begin{table}[b]
    \caption{This is an example table.}    
    \centering
    \begin{tabular}{ccc}
            \hline
            Column 1 & Column 2 & Column 3 \\  
            \hline
            Cell 1 & Cell 2 & Cell 3\\ 
            Cell 4 & Cell 5 & Cell 6 \\
            \hline
            \end{tabular}

    \label{tab:1}
\end{table}

While they should appear embedded in the text at place of first citation in initial submission, we ask that Figures and Tables also be submitted as separate, raw files should your submission move past the revision stage. This will make for easier production should the paper be accepted. At that stage, figures should be supplied as Adobe Portable Document Format (PDF), PostScript (PS), or Encapsulated PostScript (EPS) for illustrations or diagrams; Tagged Image File Format (TIFF), JPEG, PNG, PhotoShop (PSD), EPS, or PDF for photography or microscopy.

Figure legends (i.e., captions) should be listed immediately after each figure as it appears embedded in the text. Each figure legend should start with a short boldface title beginning with (for example) \textbf{Figure 1}. See Figure \ref{fig:1} for example. All figure titles should be phrases or sentences; do not mix the two styles. No single legend should be longer than about 200 words. Nomenclature, abbreviations, symbols, and units used in a figure legend (and in the figure itself) should match those used in the text. Use roman letters (A, B, C, etc.) to label individual figure panels when first mentioned in a legend: \textbf{(A)}, \textbf{(B)}, etc.

\begin{figure}[h]
    \centering
    \begin{subfigure}{0.4\textwidth}
        \includegraphics[width=0.9\textwidth, height=2in]{fig 1}
        \caption{\label{fig:2a}}
    \end{subfigure}
    \begin{subfigure}{0.4\textwidth}
        \includegraphics[width=0.9\textwidth, height=2in]{fig 2}
        \caption{\label{fig:2b}}
    \end{subfigure}
    \caption{This is an example of a figure consisting of multiple panels.     (\subref{fig:2a}) This is the first panel. (\subref{fig:2b}) This is the second panel.}
    \label{fig:2}
\end{figure}

%%%%%% Main Text %%%%%%

\section{Introduction}
The manuscript should start with a brief introduction that lays out the problem addressed by the research and describes the paper’s importance. The scientific question being investigated should be described in detail. The introduction should provide sufficient background information to make the article understandable to readers in other disciplines and provide enough context to ensure that the implications of the experimental findings are clear. 

\section{Methods}
The methods section should provide sufficient information to allow replication of the results.

\subsection{Experimental Design}
Begin with a section titled Experimental Design describing the objectives and design of the study as well as prespecified components. In this section, please include a diagram or flowchart to show the entire experimental design and illustrate the most significant elements: materials, treatments, measurements, data collection, methods of data analysis. This will facilitate the editors, reviewers and readers to understand and follow the whole concept, design, and results.

\subsection{Statistical Analysis}
Include a section titled Statistical Analysis that fully describes the statistical methods with enough detail to enable a knowledgeable reader with access to the original data to verify the results. The values for \textit{N}, \textit{P}, and the specific statistical test performed for each experiment should be included in the appropriate figure legend or main text. 

This section should be broken up by subheadings. Under exceptional circumstances, when a particularly lengthy description is required, a portion of the materials and methods can be included in the Supplementary Materials. 

\section{Results}
The results should describe the experiments performed and the findings observed. The results section should be divided into subsections to delineate different experimental themes. Subheadings should either be all phrases or all complete sentences. All data must be shown either in the main text or in the Supplementary Materials.
\begin{itemize}
    \item All data should be presented in the Results. No data should be presented for the first time in the Discussion. Data (such as from Western blots) should be appropriately quantified.
    \item Subheadings must be either all complete sentences or all phrases. They should be brief, ideally less than 10 words. Subheadings should not end in a period. Your paper may have as many subheadings as are necessary.
    \item Figures and tables must be called out in numerical order. For example, the first mention of any panel of Figure 3 cannot precede the first mention of all panels of Figure 2. The supplementary figures (for example, Figure S1) and tables (Table S1) must also be called out in numerical order.
\end{itemize}

\section{Discussion}
Include a Discussion that summarizes (but does not merely repeat) your conclusions, elaborates on their implications or significance, and compares with previously published results. There should be a paragraph outlining the limitations of your results and interpretation, as well as a discussion of the steps that need to be taken for the findings to be applied. Please avoid claims of priority. 

\begin{flushleft}
\textbf{Acknowledgments and other statements}: All acknowledgments statements should be included at the very end of the manuscript before the references. The following sections, when needed, should be listed in the order below and include the subheads in bold followed by single paragraphs including the information specified. 
\end{flushleft}

\subsection*{Ethical Approval} 
For all animal experimentation described in the manuscript, authors must state their adherence to the NIH Guide for the Care and Use of Laboratory Animals or the equivalent. Informed consent must be obtained for studies on humans after the nature and possible consequences of the studies are explained. A statement that informed consent was obtained must also appear in the manuscript. All research on humans must have approval from the institutional IRB (Institutional Review Board) or an equivalent body. 

\subsection*{Data Availability}
Submission of a manuscript to \textit{Health Data Science} implies that the data is freely available upon request or has deposited to an open database, like NCBI. If data are in an archive, include the accession number or a placeholder for it. Also include any materials that must be obtained through an MTA.  

\subsection*{Funding}
Authors must state how the research and publication of their article was funded, by naming financially supporting body(s) (written out in full) followed by associated grant number(s) in square brackets (if applicable), for example: “This work was supported by the Engineering and Physical Sciences Research Council [grant numbers xxxx, yyyy]; the National Science Foundation [grant number zzzz]; and a Leverhulme Trust Research Project Grant”. If the research did not receive specific funding but was performed as part of the employment of the authors, please name this employer. If the funder was involved in the manuscript writing, editing, approval, or decision to publish, please declare this.

\subsection*{Authors' Contributions} 
Describe contributions of each author (use initials) to the paper. 

\medskip Examples:

``S. Z. conceived the idea and designed the experiments.''

``E. F. M. and J. F. S. conducted the experiments.''

``All authors contributed equally to the writing of the manuscript.''

\subsection*{Conflicts of Interest}
Conflicts of interest (COIs, also known as ``competing interests'') occur when issues outside research could be reasonably perceived to affect the neutrality or objectivity of the work or its assessment. For more information, see our publication ethics policy. Authors must declare all potential interests – whether or not they actually had an influence – in a ‘Conflicts of Interest’ section in the Acknowledgements, which should explain why the interest may be a conflict. If there are none, the authors should state ``The author(s) declare(s) that there is no conflict of interest regarding the publication of this article.'' Submitting authors are responsible for coauthors declaring their interests. Declared conflicts of interest will be considered by the editor and reviewers and included in the published article.

The involvement of anyone other than the authors who 1) has an interest in the outcome of the work; 2) is affiliated to an organization with such an interest; or 3) was employed or paid by a funder, in the commissioning, conception, planning, design, conduct, or analysis of the work, the preparation or editing of the manuscript, or the decision to publish must be declared.

\subsection*{Acknowledgments}
Thank others for any contributions. Anyone who made a contribution to the research or manuscript, but who is not a listed author, should be acknowledged (with their permission).

\section*{Supplementary Materials}
Include supporting text (including notes, supplementary materials and methods, tables, and figures) at the end of the main manuscript file, in a separate section titled Supplementary Materials, if this can be easily done and if the total file size does not exceed 6 MB. Alternatively, Supplementary Materials can be included as a separate file that can be uploaded as the final figure file within the 6 MB upload limit. 

Supplementary figures should be embedded in the file in order, with the legends directly below the figure. 

Any references cited in the Supplementary Materials must already appear in the reference list; no separate supplementary reference list should be created.
 
Supplementary Materials may include additional author notes—for example, a list of group authors.

If you have any Supplemental Materials please list them by sections in the following order: supplementary materials and methods (if any), supplementary figures, supplementary tables, and other supplementary files (such as movies, data, interactive images, or database files). Be sure to submit all Supplementary Materials with the manuscript. Materials should be named as follows: 

\medskip Figures S1 to S\#

Tables S1 to S\#

Movies S1 to S\#

Data files S1 to S\#

Audio files S1 to S\#

\medskip

\section*{References}
There is only one reference list for all sources cited in the main text, figure and table legends, and Supplementary Materials. Do not include a second reference list in the Supplementary Materials section. Include references cited only in the Supplementary Materials at the end of the reference section of the main text; reference numbering should continue as if the Supplementary Materials are a continuation of the main text. References cited only in the Supplementary Materials section are not counted toward length guidelines.

Authors are responsible for ensuring that the information in each reference is complete and accurate. All data must be cited and references to “data not shown” or citations to unpublished results are permitted.

Please do not include any extraneous language such as explanatory notes as part of a reference to a given source. \textit{Health Data Science} prefers that manuscripts do not include end notes; if information is important enough to include, please put into main text.  If you need to include notes, please explain why they are needed in your cover letter to the editor.

\printbibliography

\end{document}
